\section{Use Case}
\xiangyu{We should give a summary sentence in the beginning of this section. (e.g.,\sys{} is not limited to a single application type but is designed to address a diverse range of use cases. We discuss three representative ones in this section.)}

\subsection{GPU Memory Tracer}\label{sec:use_case_gpu_memory_tracer}
\noindent \sys instruments LD/ST sites via PTX trampolines and logs address, size, warp/SM IDs, and timestamps into a shared eBPF map. A user-space consumer derives bandwidth, coalescing, and contention metrics with sub-\(\mu\)s latency, enabling quick diagnosis of bank conflicts or uncoalesced accesses without relaunching kernels.

\subsection{Fair Share Scheduling}\label{sec:use_case_fair_share}
\noindent Block-level entry/exit probes report per-block runtimes into shared maps. A host scheduler adjusts future block quotas to equalize progress across tenants, yielding near-ideal Jain's fairness with \(<2\%\) overhead while remaining transparent to application code and CUDA streams.

\subsection{ML Green: CPU–GPU Collaborative Inference}
\noindent Modern Large Language Models often exceed the capacity of on-chip GPU memory\xiangyu{Why do we need to mention energy generation timing here? If it is necessary, we should separate it into 2 sentences. The first sentence says that LLM requires more memory than on-chip GPU limit; the second sentence mentions that the energy generation time varies because of the sunlight and wind is unstable; the third maybe say that this cause some miss match between energy demand in LLM and the energy supply from various energy generation sources}. For models that exceed on-device memory or under renewable-aware power budgets, \sys traces tensor access patterns to guide split execution akin to Neuro-Surgeon~\cite{neurosurgeon}.

\xiangyu{The central challenge is three-fold --> Realizing this goal faces three typical challenges}: determining which model state should remain in GPU memory versus CPU RAM; prefetching evicted data just in time; and shifting computation to CPU when GPU energy consumption exceeds available renewable power (e.g., at night without solar).

\sys{} probes memory accesses to reveal when and how often each tensor is accessed. Armed with these traces, engineers can emulate Neuro-Surgeon's partitioning logic: high-priority layers stay in GPU memory while lower-priority elements live in CPU memory until needed.

We compare three \xiangyu{typical} eviction strategies under this CPU–GPU split:
\begin{itemize}[leftmargin=*,itemsep=0pt]
\item \textbf{LRU (Least Recently Used):} evict the tensor that has gone longest without use;
\item \textbf{LFU (Least Frequently Used):} retain the most–accessed tensors, which can outperform LRU when a small subset is hot;
\item \textbf{Full CPU, Full GPU offloading:} dynamically stream each tensor on demand (Full CPU), or inject the entire model into GPU memory (Full GPU).
\end{itemize}

Policies such as LRU/LFU and full-CPU/full-GPU offload are evaluated online; access and energy traces steer eviction and prefetch so transfers overlap compute under varying power availability. \sys's fine-grained energy and caching logs help tune the cache to hide transfer overhead behind GPU compute under renewable energy constraints.

